% ============================================================
% masterclass-40page-writer · 可直接复制的 LaTeX 前置模板
% 已在 tectonic + ctexart + 中文 环境实测通过（多篇 40+ 页文章验证）
% 用法：把本文件 \begin{document} 之前的内容整段复制为新主文件的 preamble
% ============================================================

\documentclass[11pt,a4paper]{ctexart}
\usepackage[utf8]{inputenc}
\usepackage{geometry}
\geometry{top=2.2cm,bottom=2.2cm,left=2.5cm,right=2.5cm}
\usepackage{amsmath,amssymb,amsthm}
\usepackage{graphicx}
\usepackage{xcolor}
\usepackage{tikz}
\usetikzlibrary{arrows.meta,positioning,calc,shapes.geometric,patterns,decorations.pathreplacing,fit}
\usepackage{listings}
\usepackage{booktabs}
\usepackage{tcolorbox}
\tcbuselibrary{skins,breakable}
\usepackage{algorithm}
\usepackage{algpseudocode}
\usepackage{hyperref}
\hypersetup{
    colorlinks=true,
    linkcolor=blue!70!black,
    citecolor=green!60!black,
    urlcolor=blue!70!black
}

\definecolor{codegreen}{rgb}{0,0.6,0}
\definecolor{codegray}{rgb}{0.5,0.5,0.5}
\definecolor{codepurple}{rgb}{0.58,0,0.82}
\definecolor{backcolour}{rgb}{0.97,0.97,0.97}

\lstset{
    backgroundcolor=\color{backcolour},
    commentstyle=\color{codegreen},
    keywordstyle=\color{magenta},
    numberstyle=\tiny\color{codegray},
    stringstyle=\color{codepurple},
    basicstyle=\ttfamily\small,
    breakatwhitespace=false,
    breaklines=true,
    captionpos=b,
    keepspaces=true,
    numbers=left,
    numbersep=5pt,
    showspaces=false,
    showstringspaces=false,
    showtabs=false,
    tabsize=2,
    frame=single,
    rulecolor=\color{black!20}
}

% ---- 六种教学盒子：高中生友好长文的核心装置 ----
% 1) 卷首引语
\newtcolorbox{quotebox}{
    colback=gray!5!white, colframe=gray!40!black, fonttitle=\bfseries,
    arc=2mm, boxrule=0.8pt, left=12pt, right=12pt, top=8pt, bottom=8pt
}
% 2) 生活类比（橙色）—— 每个元概念必须先从这里切入
\newtcolorbox{tipbox}[1]{
    colback=orange!4!white, colframe=orange!70!black, fonttitle=\bfseries,
    title={#1}, arc=2mm, boxrule=0.8pt,
    left=10pt, right=10pt, top=6pt, bottom=6pt
}
% 3) 想一想（黄色）—— 停下来让读者思考
\newtcolorbox{thinkbox}{
    colback=yellow!6!white, colframe=yellow!50!black, fonttitle=\bfseries,
    title={想一想}, arc=2mm, boxrule=0.8pt,
    left=10pt, right=10pt, top=6pt, bottom=6pt
}
% 4) 练习 / 手算（青色）—— 必须给答案
\newtcolorbox{exbox}[1]{
    colback=teal!4!white, colframe=teal!60!black, fonttitle=\bfseries,
    title={#1}, arc=2mm, boxrule=0.8pt,
    left=10pt, right=10pt, top=6pt, bottom=6pt
}
% 5) 本篇小结（蓝色）—— 每篇末尾必用，标题写「本篇小结（你已会什么）」
\newtcolorbox{stepbox}[1]{
    colback=blue!3!white, colframe=blue!60!black, fonttitle=\bfseries,
    title={#1}, arc=2mm, boxrule=0.8pt,
    left=10pt, right=10pt, top=6pt, bottom=6pt
}
% 6) 进阶数学（紫色）—— 明确标注「第一次读可跳过」
\newtcolorbox{rigorbox}[1]{
    colback=purple!3!white, colframe=purple!70!black, fonttitle=\bfseries,
    title={#1}, arc=2mm, boxrule=0.8pt,
    left=10pt, right=10pt, top=6pt, bottom=6pt
}
% 7) 风险警示（红色）—— 【补齐】SKILL.md 列出了 warnbox，但旧版 preamble 漏定义了，
%    直接 \begin{warnbox}{...} 会报 Undefined control sequence 并导致 XeTeX 崩溃。
\newtcolorbox{warnbox}[1]{
    colback=red!4!white, colframe=red!70!black, fonttitle=\bfseries,
    title={#1}, arc=2mm, boxrule=0.8pt,
    left=10pt, right=10pt, top=6pt, bottom=6pt
}

\setcounter{tocdepth}{2}   % 40 页长文必须有目录

\title{\textbf{\LARGE <主标题>}\\[0.5em]
\textbf{\Large <副标题>}\\[0.5em]
{\large —— 一本给高中生也能读懂的 40 页通识长文}}
\author{\large 散步}
\date{}

\begin{document}
\maketitle
% 卷首引语（两句真实可考的名言/论文原话）
% \begin{quotebox} ... \end{quotebox}
\tableofcontents
\newpage
% \input{<parts 目录>/p1_xxx}
\end{document}
